\IEEEPARstart{G}{lobal} \add{navigation satellite systems are unreliable indoors, where attenuation, blockage and multipath degrade their signals. Radio-frequency alternatives, such as Wi-Fi, Bluetooth Low Energy, ultra-wideband and 5G positioning~\cite{Faragher2015,Alarifi2016,Zafari2019,Italiano2025}, are widely deployed, but in cluttered interiors multipath and interference typically limit them to meter-level accuracy unless dense infrastructure or dedicated calibration is used~\cite{Lymberopoulos2015,Leitch2023}.}

\add{Optical wireless positioning, and visible light positioning (VLP) in particular, exploits the LED lighting already present in buildings~\cite{Wang2024b,Zhu2025}. Light does not cross walls, is immune to radio interference and propagates mostly along a deterministic line of sight, so that centimeter-level accuracy can be obtained with inexpensive components~\cite{Trevlakis2023}, and the same luminaires can also carry optical wireless communication~\cite{Haas:16}. Among the possible observables, the received signal strength (RSS) of a photodiode (PD) is the simplest, since it requires neither synchronization nor imaging~\cite{Zhuang2018}.}

\add{Most VLP systems follow a multi-LED paradigm, in which several luminaires of known position act as anchors and the receiver trilaterates its position~\cite{Zhuang2018,Wang2024b}. Three or four LEDs must then be visible simultaneously to resolve a three-dimensional (3D) position~\cite{Plets2019}, a condition that sparse lighting, corridors or obstacles often violate, and every lamp needs its own identifier or modulation, which raises the cost of deployment and maintenance~\cite{Cossu2022}. Positioning from a single LED therefore removes a major barrier to adoption~\cite{Rekkas2023}.}

\add{With a single LED, one RSS measurement carries too little information, and the receiver must supply the missing diversity. Most single-LED systems do so with several PDs of different orientations. With one horizontal and three tilted PDs, the ratios of their RSS determine the 3D position in closed form, with a mean error of 2.52~cm in simulation~\cite{Qin2020}, and the same receiver reaches 9.27~cm in experiments with Bayesian fingerprint matching~\cite{Qin2021}. Learning-based estimators, from recurrent networks~\cite{Chen2021} and extreme learning machines~\cite{Zhang2024} to deep residual shrinkage networks~\cite{Li2024}, report errors between 0.9 and 5.9~cm with three or four PDs, and a four-PD receiver combined with orientation sensing handles arbitrary orientations with a mean error of 1.77~cm~\cite{Ma2024}; least squares on a four-PD receiver gives 7~cm at the 90th percentile~\cite{Zia-Ul-Mustafa2023}. Such receivers observe all directions at once, but they multiply the detectors and their acquisition channels, whose gains and responses must be matched or calibrated.}

\add{A single PD can create the same diversity over time by being reoriented. On a rotating arm, it has been combined with an extreme learning machine that accounts for first-order reflections~\cite{Liu2022}, and with least squares and a Cram\'er--Rao bound that selects the arm length and the rotation angles, reaching 4.96~cm in experiments with four rotations~\cite{Shi2025}; the bound under signal-dependent shot noise has also been used to compare such receivers and to adapt their tilt~\cite{Shi2026}. A PD tilted on a rotating base has been analyzed with a bound that guides the choice of the tilt and positioned with a least-squares-initialized grey wolf optimizer, at the centimeter level in a 0.5$\times$0.5$\times$0.3~m setup~\cite{Wang2024}. The diversity can also be moved to the infrastructure, by reorienting the source~\cite{Chassagne2025,Acuna2026} or with intelligent reflecting surfaces~\cite{ShiJin2025}, at the price of steerable or additional equipment around the LED.}

\add{These works establish that a single LED suffices, but several aspects limit their transfer to practice. Most of them estimate the position in two dimensions at a known height~\cite{Qin2021,Liu2022,Wang2024,Shi2025}, and many rely on fingerprints or learned models that must be retrained when the hardware or the environment changes~\cite{Rekkas2023,Glass2021}, whereas model-based estimators remain interpretable and need little data~\cite{Deng2023}. The receiver is nearly always modeled with a cosine response, although the response of commercial PDs departs from it and the mismatch degrades RSS positioning~\cite{Bastiaens2020}; a reoriented receiver is even more exposed, since it infers the direction from how the response changes between orientations. Finally, the evaluations are mostly simulations or small setups with a few tens of positions, and the effect of errors in the orientations, which a reoriented receiver never knows exactly, has received little attention.}

\add{This paper addresses these aspects with a single fixed LED and a single PD reoriented to $K$ known orientations, from whose RSS all three coordinates are estimated. The shape of the $K$ RSS over the orientations gives the direction to the LED, and their common amplitude gives the distance through a calibrated model of the LED, so that no additional measurement, detector or modification of the lighting is required. The receiver can be as lean as one tilted PD: a drone hovering under a lamp, for instance, could obtain the orientations by yawing, with its attitude given by its inertial sensors. The main contributions are the following:}
\begin{itemize}
 \item \add{A measurement model and bound for an arbitrary receiver response and arbitrary known orientations: Fisher information, identifiability condition, effect of independent and common orientation errors, and a bound-based design of the receiver tilt.}
 \item \add{Estimators for a general response, including nonlinear ratio estimators for responses that admit no linearization, such as the SQ model of commercial PDs, with a closed-form distance that yields the maximum-likelihood position, and an RSS-shape test that calibrates the receiver orientation and monitors its integrity without ground truth.}
 \item \add{An experimental validation on an automated room-scale testbed with a ceiling LED at 3~m: the mean 3D error is 3.66~cm over 125 positions at five heights, and 4.10~cm over 131 random positions estimated live with five orientations and a calibration fixed beforehand; the evaluation is repeated with three other LEDs.}
 \item \add{An analysis of what limits this accuracy and of what the system requires: measured noise against the bound, sensitivity to orientation errors, a calibration with nine points, an observation window of 10~ms and a few milliseconds of computation per position.}
\end{itemize}

\add{Section~\ref{sec:model} presents the system and measurement model, Section~\ref{sec:bounds} the identifiability and the bound, and Section~\ref{sec:estimation} the estimators and the calibration. Section~\ref{sec:testbed} describes the testbed and the datasets. Sections~\ref{sec:results} and~\ref{sec:analysis} report the accuracy and the error analysis, Section~\ref{sec:discussion} discusses the results and their limitations, and Section~\ref{sec:conclusion} concludes the paper.}
